\section{The relation}\label{SectRelation}
We show in this section that the stated skein algebra of a surface with a single boundary edge is isomorphic to $A_\Sigma$. We consider the algebra $A_\Sigma$ from last section with $\mathcal{V} = \Oq\text{-}{\rm comod}^{\rm fin}$ and $\mathcal{E} = \Oq\text{-}{\rm comod}$ at generic $q$, and prove that $\mathscr{S}(\Sigma)\simeq A_\Sigma$ as $\Oq$-comodule algebras. Actually since $A_\Sigma$ is only defined up to canonical isomorphism one may take an equality here, so we prove that $\mathscr{S}(\Sigma)$ satisfies the defining properties of $A_\Sigma$, namely that it is the internal endomorphism algebra of the empty set in ${\rm SK}_\mathcal{V}(\Sigma)$ with respect to the $\mathcal{E}$-module structure. This result is not new and was stated in a weaker form in \cite[Theorem~4.4]{LeYu}, \cite[Theorem~9.1]{LeSikora} and \cite[Remark 2.21]{Gunningham}, namely as algebras in $\rm Vect$. However, in these references one considers right $\mathcal{E}$-actions and therefore the internal skein algebra is isomorphic to the braided opposite of the stated skein algebra. The full result can still be recovered using \cite[Theorem~5.3]{Faitg} or \cite{Korinman}, which give an isomorphism between the opposite of the stated skein algebras and Alekseev--Grosse--Schomerus-algebras, which are themselves isomorphic to internal skein algebras by \cite{BBJ}. Our approach here is more direct and uses only skein theory.

We need a natural isomorphism $\St_W\colon \Hom_{{\rm SK}_\mathcal{V}(\Sigma)}(W \rhd \varnothing,\varnothing) \tilde{\Rightarrow} \Hom_{\mathcal{E}}(W, \mathscr{S}(\Sigma))$, for $W \in \mathcal{E} = \Oq\text{-}{\rm comod}$. In the case where $W = V^{\otimes n}\in {\rm TL}$ is a tensor product of standard corepresentations, and the element on the left hand side is a tangle $\alpha$ with $n$ boundary points, we want a morphism $V^{\otimes n}\to \mathscr{S}(\Sigma)$ associated to it. This is done by setting the entries, elements of $V^{\otimes n}$, as states of the tangle $\alpha$. Recall that $V$ has basis $v_+$, $v_-$ and we identify the states + and $-$ with these elements. As~in~Remark~\ref{rmkFormActOnSS}, we allow formal linear combination of states for stated tangles. In this context, the defining relations of stated skein algebras correspond exactly to naturality conditions, see~Proposition~\ref{propStNat}.
This gives the idea of how to deal with the objects of the full subcategory ${\rm TL} \subseteq \mathcal{V}$ of objects of the form $V^{\otimes n}$, and this extends to $\mathcal{E} \simeq \Free({\rm TL})$ by Proposition~\ref{propIntHomFree}. Note that one still has an action $\rhd \colon {\rm TL} \otimes {\rm Sk}_{{\rm TL}}(\Sigma) \to {\rm Sk}_{{\rm TL}}(\Sigma)$ which is the restriction of $\rhd \colon \mathcal{V} \otimes {\rm Sk}_\mathcal{V}(\Sigma) \to {\rm Sk}_\mathcal{V}(\Sigma)$.
\begin{Definition}
Let $W = V^{\otimes n}\in {\rm TL}$ and $\alpha \in \Hom_{{\rm Sk}_{{\rm TL}}(\Sigma)}(W \rhd \varnothing,\varnothing)$. By Theorem~\ref{thmSkAsTangles}, $\alpha$~can be represented by a (linear combination of) tangle(s), still denoted by $\alpha$, with $n$ ordered boundary points near the boundary edge, which is well defined up to isotopy and Kauffman-bracket relations.

Graphically, we set
\[
\St_W\Bigg(
\begin{tikzpicture}[scale = 1, baseline = 8pt]
\node [rightymissy, minimum width = 1cm] (A) at (1.5,1) {$\alpha$};
\draw (0,0) -- ++(2,0) node[pos = 0,below left =-2pt]{\footnotesize $\partial \Sigma$} node[pos = 0.5,below]{\footnotesize $\Sigma$};
\draw (0,0) -- ++ (0,1) node[pos = 0.5,left]{\footnotesize $[0,1]$};
\draw[blue] (0.5,0) node[above left = -2pt]{\tiny 1} .. controls (0.5,1) .. (A);
\draw[olive] (1,0) node[above left = -2pt]{\tiny $\cdots$} .. controls (1,0.7) .. (A);
\draw[red] (1.5,0) node[above left = -2pt]{\tiny $n$} -- (A);
\end{tikzpicture} \Bigg) \ (v_{\varepsilon_1}\otimes \cdots \otimes v_{\varepsilon_n}) := \begin{tikzpicture}[scale = 1, baseline = 8pt]
\node [rightymissy, minimum width = 1cm] (A) at (1.5,1) {$\alpha$};
\draw (0,0) -- ++(2,0) ;
\draw[->] (0,0) -- ++ (0,1);
\draw[blue] (0,0.75) node[ left = -2pt]{\footnotesize $\varepsilon_1$} -- (A);
\draw[olive] (0,0.5) node[ above , rotate = 90]{\tiny $\cdots$} .. controls (1,0.5) .. (A);
\draw[red] (0,0.25) node[ left = -2pt]{\footnotesize $\varepsilon_n$} .. controls (1.5,0.25) .. (A);
\end{tikzpicture}\,,\qquad \varepsilon_i \in \{\pm\}.
\]
For brevity we denote this last stated tangle by $_{\vec{\varepsilon}}\,\alpha$ and $v_{\vec{\varepsilon}} = v_{\varepsilon_1}\otimes \cdots \otimes v_{\varepsilon_n}$. Note that in this figure there are two implicit sums, $\alpha$ is a linear combination of tangles and an element of $V^{\otimes n}$ is a linear combination of $v_{\vec{\varepsilon}}$'s. They will remain implicit in the following.

In the context of stated skein algebras one needs the boundary points of the tangle $\alpha$ to be above the boundary arc though they are on the bottom at its right in the context of morphisms in ${\rm Sk}_{{\rm TL}}(\Sigma)$. One wants to simply push the boundary points through the boundary edge and up above it, but without braiding the strands with one another. Therefore we use a global automorphism of $\Sigma \times [0,1]$. Consider an isotopy of the identity on $\Sigma \times [0,1]$ which is trivial far from the corner and in it pushes the bottom boundary through and up above the edge. It results in the global automorphism $\psi_{hv}$ of $\Sigma \times [0,1]$ which maps a tangle ``from skein categories'' to one ``from stated skein algebras'' and preserves the order as desired, namely the bottom points on the leftmost will end up above. We will still denote the modified tangle $\psi_{hv}(\alpha)$ by $\alpha$. Now, it only needs states to give an element of $\mathscr{S}(\Sigma)$.

We set
\[
\St_W(\alpha) := \begin{cases}
V^{\otimes n} \to \mathscr{S}(\Sigma),
\\
v_{\vec{\varepsilon}} \mapsto {}_{\vec{\varepsilon}}\,\alpha,
\end{cases}
\]
so where ${}_{\vec{\varepsilon}}\,\alpha$ is the tangle $\psi_{hv}(\alpha)$ with states $\varepsilon_1,\dots,\varepsilon_n$ from top to bottom. It is well defined because the tangle representing $\alpha$ is well defined up to isotopy and Kauffman-bracket relations, which are quotiented out in $\mathscr{S}(\Sigma)$.
\end{Definition}

\begin{Proposition}\label{propStOqmorph}
Given a tangle $\alpha$, the map $\St_W(\alpha)\colon V^{\otimes n}\to\mathscr{S}(\Sigma)$ is an $\Oq$-comodule morphism.
\end{Proposition}
\begin{proof}
Note that we still see $\mathscr{S}(\Sigma)$ as a right $\Oq$-comodule even though we draw the edge at the left. Its comodule structure is given by $\Delta(_{\vec{\varepsilon}}\,\alpha) = \sum_{\vec{\eta} \in \{\pm\}^n} {}_{\vec{\eta}}\alpha \otimes {}_{\vec{\eta}}\beta_{\vec{\varepsilon}}$, and here~${}_{\vec{\eta}}\beta_{\vec{\varepsilon}}$ is a~product $\prod_i {}_{\eta_i}\beta_{\varepsilon_i}$ in $\mathscr{S}(B)$. The comodule structure on $V^{\otimes n}$ is given by $\Delta(v_{\vec{\varepsilon}}) = \sum_{\vec{\eta}} v_{\vec{\eta}} \otimes \prod_i x_{\eta_i,\varepsilon_i}$, where $x_{\eta,\varepsilon}$~is the $v_\eta$ part of $\Delta(v_\varepsilon)$. Namely, $x_{+,+} = a$, $x_{+,-}=b$, $x_{-,+}=c$ and $x_{-,-}=d$. Under the isomorphism $\Oq \simeq \mathscr{S}(B)$, $x_{\eta,\varepsilon}\mapsto {}_\eta\beta_{\varepsilon}$.
 Thus $\Delta(\St_W(\alpha)(v_{\vec{\varepsilon}})) = \Delta(_{\vec{\varepsilon}}\,\alpha) = \sum_{\vec{\eta}} \St_W(\alpha)(v_{\vec{\eta}}) \otimes {}_{\vec{\eta}}\beta_{\vec{\varepsilon}} = (\St_W(\alpha)\otimes {\rm Id}_{\mathscr{S}(B)})(\Delta(v_{\vec{\varepsilon}}))$.
\end{proof}

\begin{Proposition}\label{propStNat}
The maps $\St_W \colon \Hom_{{\rm Sk}_{{\rm TL}}(\Sigma)}(W\rhd \varnothing,\varnothing) \to \Hom_{\Oq\text{-}{\rm comod}}(W, \mathscr{S}(\Sigma))$, $W \in {\rm TL}$, define a natural transformation $\St\colon \Hom_{{\rm Sk}_{{\rm TL}}(\Sigma)}(-\rhd \varnothing,\varnothing) \Rightarrow \Hom_{\mathcal{E}}(-, \mathscr{S}(\Sigma))$ between functors ${\rm TL}^{\rm op}\to {\rm Vect}_k$.
\end{Proposition}
\begin{proof}
For $g \in \Hom_{{\rm TL}}(V^{\otimes n},V^{\otimes m})$, $\alpha \in \Hom_{{\rm Sk}_{{\rm TL}}(\Sigma)}(V^{\otimes m}\rhd \varnothing,\varnothing)$ and $v_{\vec{\varepsilon}} \in V^{\otimes n}$
, one needs to check that $\St_{V^{\otimes m}}(\alpha)(g(v_{\vec{\varepsilon}})) = \St_{V^{\otimes n}}(\alpha \circ (g\rhd {\rm Id}_\varnothing))(v_{\vec{\varepsilon}})$. Morphisms of ${\rm TL}$ are generated under composition and juxtaposition by identities, caps and cups, so one only needs to prove the result for these last two.
We will need some explicit computations of these caps and cups. Recall from Remark \ref{rmkUnorTanOq}, or \cite[Theorem~4.2]{Tingley}, that to do so one uses the isomorphism
\[
\varphi \colon\ \begin{cases}
V \to V^*,
\\
v_+ \mapsto -q^{\frac{5}{2}} v_-^*,
\\ v_- \mapsto q^{\frac{1}{2}}v_+^*\end{cases}
\]
 of Definition~\ref{defOqcomods} and sets $\cap = \lcap \circ (\varphi \otimes {\rm Id}_V)$ and $\cup = ({\rm Id}_V \otimes \varphi^{-1})\circ \lcup$, where $\lcap$ and $\lcup$ are the usual ${\rm ev}$ and ${\rm coev}$ in ${\rm Vect}_k^{\rm fin}$.

Let $g =\ $\begin{tikzpicture}[xscale = 0.18, yscale = 0.25, baseline = 0pt]
\draw (0,0) -- ++(0,1) ;
\draw (1,0) -- ++(0,1) ;
\draw (2,0) -- ++(0,1) ;
\draw (3,0) arc(180:0:0.5 and 0.75);
\draw (5,0) -- ++(0,1) ;
\draw (6,0)-- ++(0,1) ;
\draw (7,0) -- ++(0,1) ;
\end{tikzpicture}$\ = {\rm Id}_V^{\otimes k} \otimes \cap \otimes {\rm Id}_V^{\otimes n-k}\colon V^{\otimes n+2} \to V^{\otimes n}$, $\alpha \in \Hom_{{\rm Sk}_{{\rm TL}}(\Sigma)}(V^{\otimes n}\rhd \varnothing,\varnothing)$ and $v=(v_{\varepsilon_1}\otimes \cdots \otimes v_{\varepsilon_k})\otimes v_{\mu}\otimes v_{\nu} \otimes (v_{\varepsilon_{k+1}}\otimes \cdots\otimes v_{\varepsilon_{n}}) \in V^{\otimes n+2}$. We want to compare
\begin{gather*}
\St_{V^{\otimes n}}(\alpha)(g(v)) = \St_{V^{\otimes n}}(\alpha)(v_{\vec{\varepsilon}}\ .\cap(v_{\mu}\otimes v_{\nu})) =\begin{tikzpicture}[scale = 1, baseline = 8pt]
\node [rightymissy, minimum width = 1cm] (A) at (1.5,1) {$\alpha$};
\draw (0,0) -- ++(2,0) ;
\draw (0,0) -- ++ (0,1);
\draw (0,0.75) node[ left = -2pt]{\footnotesize $\varepsilon_1$} -- (A);
\draw (0,0.5) node[ above , rotate = 90]{\tiny $\cdots$} .. controls (1,0.5) .. (A);
\draw (0,0.25) node[ left = -2pt]{\footnotesize $\varepsilon_n$} .. controls (1.5,0.25) .. (A);
\end{tikzpicture}\ .\cap(v_{\mu}\otimes v_{\nu})
\end{gather*}
and
\begin{gather*}
\St_{V^{\otimes n+2}}(\alpha \circ (g\rhd {\rm Id}_\varnothing))(v)=\St_{V^{\otimes n+2}}\Bigg(\ \begin{tikzpicture}[scale = 0.7, baseline = 8pt]
\node [rightymissy, minimum width = 1cm] (A) at (2.5,1.5) {$\alpha$};
\draw (0,0) -- ++(3,0);
\draw (0,0) -- ++ (0,1.5);
\draw(0.5,0) .. controls (0.5,1.5) .. (A);
\draw(1,0) .. controls (1,1.2) .. (A);
\draw (1.5,0) arc(180:0:0.25 and 0.5);
\draw(2.5,0) -- (A);
\end{tikzpicture}
\Bigg)(v) = \begin{tikzpicture}[scale = 0.8, baseline = 8pt]
\node [rightymissy, minimum width = 1cm] (A) at (2.5,1.5) {$\alpha$};
\draw (0,0) -- ++(3,0) ;
\draw (0,0) -- ++ (0,1.5);
\draw (0,1.25) node[ left = -2pt]{\footnotesize $\varepsilon_1$} -- (A);
\draw (0,1) node[ above , rotate = 90]{\tiny $\cdots$} .. controls (1,1) .. (A);
\draw (0,0.75) node[ left = -2pt]{\footnotesize $\mu$} arc(90:-90:0.5 and 0.125) node[ left = -2pt]{\footnotesize $\nu$} ;
\draw (0,0.25) node[ left = -2pt]{\footnotesize $\varepsilon_n$} .. controls (1.5,0.25) .. (A);
\end{tikzpicture}
\\ \hphantom{\St_{V^{\otimes n+2}}(\alpha \circ (g\rhd {\rm Id}_\varnothing))(v)}
{}=\begin{tikzpicture}[scale = 1, baseline = 8pt]
\node [rightymissy, minimum width = 1cm] (A) at (1.5,1) {$\alpha$};
\draw (0,0) -- ++(2,0) ;
\draw (0,0) -- ++ (0,1);
\draw (0,0.75) node[ left = -2pt]{\footnotesize $\varepsilon_1$} -- (A);
\draw (0,0.5) node[ above , rotate = 90]{\tiny $\cdots$} .. controls (1,0.5) .. (A);
\draw (0,0.25) node[ left = -2pt]{\footnotesize $\varepsilon_n$} .. controls (1.5,0.25) .. (A);
\end{tikzpicture}. {}^{\mu}_{\nu}\text{\reflectbox{$C$}}.
\end{gather*}
One simply needs to check that the coefficient ${}^{\mu}_{\nu}\reflectbox{$C$}$ from the left boundary skein relations of Proposition~\ref{propLeftStSkRel} coincides with $\cap(v_{\mu}\otimes v_{\nu})$ and indeed $\cap(v_+ \otimes v_+) = {\rm ev}\big({-}q^{\frac{5}{2}}v_-^* \otimes v_+\big) = 0={}^+_+\reflectbox{$C$}$, $\cap(v_- \otimes v_-) = {\rm ev}\big(q^{\frac{1}{2}}v_+^* \otimes v_-\big) = 0={}^-_-\reflectbox{$C$}$, $\cap(v_+ \otimes v_-) = {\rm ev}\big({-}q^{\frac{5}{2}}v_-^* \otimes v_-\big) = -q^{\frac{5}{2}}={}^+_-\reflectbox{$C$}$ and $\cap(v_- \otimes v_+) = {\rm ev}\big(q^{\frac{1}{2}}v_+^* \otimes v_+\big) = q^{\frac{1}{2}}={}^-_+\reflectbox{$C$}$.

Now let $g =\ $\begin{tikzpicture}[xscale = 0.18, yscale = -0.25, baseline =-7pt]
\draw (0,0) -- ++(0,1) ;
\draw (1,0) -- ++(0,1) ;
\draw (2,0) -- ++(0,1) ;
\draw (3,0) arc(180:0:0.5 and 0.75);
\draw (5,0) -- ++(0,1) ;
\draw (6,0)-- ++(0,1) ;
\draw (7,0) -- ++(0,1) ;
\end{tikzpicture}$\ = {\rm Id}_V^{\otimes k} \otimes \cup \otimes {\rm Id}_V^{\otimes n-k}\colon V^{\otimes n} \to V^{\otimes n+2}$, $\alpha \in \Hom_{{\rm Sk}_{{\rm TL}}(\Sigma)}\big(V^{\otimes n+2}\rhd \varnothing,\varnothing\big)$ and~$v=v_{\varepsilon_1}\otimes \cdots \otimes v_{\varepsilon_{n}} \in V^{\otimes n}$. One can directly compute $\cup (1) = \big({\rm Id}_V \otimes \varphi^{-1}\big)\circ
\mathop{\rm coev}(1) = \big({\rm Id}_V \otimes \varphi^{-1}\big) (v_- \otimes v_-^* + v_+ \otimes v_+^*) = -q^{-\frac{5}{2}} v_-\otimes v_+ + q^{-\frac{1}{2}} v_+ \otimes v_-$.
 We want to compare
\begin{gather*}
{\rm St}_{V^{\otimes n+2}}(\alpha)(g(v)) = {\rm St}_{V^{\otimes n+2}}(\alpha)((v_{\varepsilon_1} \otimes\cdots v_{\varepsilon_k}) \otimes  \big({-}q^{-\frac{5}{2}} v_- \otimes  v_+  +  q^{-\frac{1}{2}} v_+  \otimes   v_-\big) \\
\hphantom{\St_{V^{\otimes n+2}}(\alpha)(g(v))=}
\otimes (v_{\varepsilon_{k+1}} \otimes  \cdots v_{\varepsilon_{n}}))
\\  \hphantom{\St_{V^{\otimes n+2}}(\alpha)(g(v))}
{}=-q^{-\frac{5}{2}}\begin{tikzpicture}[scale = 1, baseline = 8pt]
\node [rightymissy, minimum width = 1cm] (A) at (2.5,1.5) {$\alpha$};
\draw (0,0) -- ++(3,0) ;
\draw (0,0) -- ++ (0,1.5);
\draw (0,1.25) node[ left = -2pt]{\footnotesize $\varepsilon_1$} -- (A);
\draw (0,1) node[ above , rotate = 90]{\tiny $\cdots$} .. controls (1,1) .. (A);
\draw (0,0.75) node[ left = -2pt]{\footnotesize $-$}.. controls (1,0.75) .. (A);
\draw (0,0.5) node[ left = -2pt]{\footnotesize $+$} .. controls (1,0.5) .. (A);
\draw (0,0.25) node[ left = -2pt]{\footnotesize $\varepsilon_n$} .. controls (1.5,0.25) .. (A);
\end{tikzpicture}+q^{-\frac{1}{2}}\begin{tikzpicture}[scale = 1, baseline = 8pt]
\node [rightymissy, minimum width = 1cm] (A) at (2.5,1.5) {$\alpha$};
\draw (0,0) -- ++(3,0) ;
\draw (0,0) -- ++ (0,1.5);
\draw (0,1.25) node[ left = -2pt]{\footnotesize $\varepsilon_1$} -- (A);
\draw (0,1) node[ above , rotate = 90]{\tiny $\cdots$} .. controls (1,1) .. (A);
\draw (0,0.75) node[ left = -2pt]{\footnotesize $+$}.. controls (1,0.75) .. (A);
\draw (0,0.5) node[ left = -2pt]{\footnotesize $-$} .. controls (1,0.5) .. (A);
\draw (0,0.25) node[ left = -2pt]{\footnotesize $\varepsilon_n$} .. controls (1.5,0.25) .. (A);
\end{tikzpicture}
\end{gather*}
and
\begin{gather*}
\St_{V^{\otimes n}}(\alpha \circ (g\rhd {\rm Id}_\varnothing))(v)=\St_{V^{\otimes n}}\Bigg(\ \begin{tikzpicture}[yscale = 0.7, baseline = 8pt]
\node [rightymissy, minimum width = 1cm] (A) at (2.5,1.5) {$\alpha$};
\draw (0,0) -- ++(3,0);
\draw (0,0) -- ++ (0,1.5);
\draw(0.5,0) .. controls (0.5,1.5) .. (A);
\draw(1,0) .. controls (1,1.2) .. (A);
\draw(1.5,0.75) .. controls (1.5,1.1) .. (A);
\draw (1.5,0.75) arc(180:0:0.25 and -0.5) .. controls (2,0.9)..(A);
\draw(2.5,0) -- (A);
\end{tikzpicture}
\Bigg)(v) = \begin{tikzpicture}[scale = 1, baseline = 8pt]
\node [rightymissy, minimum width = 1cm] (A) at (2.5,1.5) {$\alpha$};
\draw (0,0) -- ++(3,0) ;
\draw (0,0) -- ++ (0,1.5);
\draw (0,1.25) node[ left = -2pt]{\footnotesize $\varepsilon_1$} -- (A);
\draw (0,1) node[ above , rotate = 90]{\tiny $\cdots$} .. controls (1,1) .. (A);
\draw (0.75,0.75) .. controls (1,0.75) .. (A);
\draw (0.75,0.75) arc(90:270:0.25 and 0.125) .. controls (1,0.5) .. (A);
\draw (0,0.25) node[ left = -2pt]{\footnotesize $\varepsilon_n$} .. controls (1.5,0.25) .. (A);
\end{tikzpicture}.
\end{gather*}
They are equal by the last left boundary skein relation of Proposition~\ref{propLeftStSkRel}.
\end{proof}

Recall from Proposition~\ref{propIntHomFree} and the discussion below that the internal endomorphism object of~$\varnothing$ with respect to the $\mathcal{E}$-action is an object $X$ equipped with an isomorphism $\Hom_{{\rm Sk}_{\rm TL}(\Sigma)}(-\rhd \varnothing,\varnothing)\allowbreak \tilde{\Rightarrow} \Hom_{\mathcal{E}}(-, X)$ in $[{\rm TL}^{\rm op},{\rm Vect}_k]$

\begin{Theorem}\label{thmRelation}
The natural transformation $\St$ is a natural isomorphism and exhibits $\mathscr{S}(\Sigma)$ as the internal endomorphism object of the empty set. Namely one can take $A_\Sigma = \mathscr{S}(\Sigma)$ as $\Oq$-comodule.
\end{Theorem}

\begin{proof}
We exhibit an inverse to $\St$. Let $W=V^{\otimes n}\in {\rm TL}$, which decomposes as a direct sum of simples $W_i$. Let $f\colon W\to \mathscr{S}(\Sigma)$ be a morphism in $\mathcal{E}$, and denote by $f_i$ its restriction on $W_i$. We~want a morphism $\St_W^{-1}(f) \in \Hom_{{\rm Sk}_{{\rm TL}}(\Sigma)}(W\rhd \varnothing,\varnothing)$, which is equivalent to a collection of morphisms ${\rm St}_W^{-1}(f_i) \in \Hom_{{\rm Sk}_{\mathcal{V}}(\Sigma)}(W_i\rhd \varnothing,\varnothing)$. Each $f_i$ is determined by its value on a single element $w_i\in W_i$ (pick a highest weight element for example), it extends to all $W_i$ by applying the $\Uq$-action (recall from Proposition~\ref{propOqcomodUqmod} that $\Oq$-comodules correspond to $\Uq$-modules). Choose any stated tangle~$a_i$ representing the element $f_i(w_i) \in \mathscr{S}(\Sigma)$. Denote by $\alpha_i$ its underlying tangle, $n_i=\abss{\partial \alpha_i}$ its number of boundary points and $\vec{\varepsilon}_i$ its states. This representative is well defined up to the boundary skein relations, the usual skein relations and isotopy. The assignment $w_i \mapsto v_{\vec{\varepsilon}_i}$ extends to a unique $ \Oq$-morphism $g_i\colon W_i \to V^{\otimes n_i}$ by applying the $\Uq$-action. We then set $\St_W^{-1}(f_i) = \alpha_i\circ (g_i \rhd {\rm Id}_\varnothing)$. Note that here~$\alpha_i$ denotes the tangle~$\alpha_i$ seen as a morphism in ${\rm Sk}_\mathcal{V}(\Sigma)$, so we actually mean $\psi_{hv}^{-1}(\alpha_i)$, the same tangle but with boundary points at the bottom instead of the left. We have to check that this definition does not depend on the representative $a_i$. Usual skein relations and isotopy do not change~$\alpha_i$ seen as a~morphism in the skein category. The boundary skein relations are equivalent to naturality using Proposition~\ref{propStNat}. Namely, another representative $a'_i$ has to be of the form $\alpha'_i= \alpha_i \circ (g \rhd {\rm Id}_\varnothing)$, for some $g \in {\rm TL}$, with states ${\vec{\varepsilon}_i}^{\ '}$ such that $g({\vec{\varepsilon}_i}^{\ '})=\vec{\varepsilon}_i$. Therefore the $g'_i$ for $a'_i$ will be such that $g_i = g \circ g'_i$. Then we simply check that $\alpha_i\circ (g_i \rhd {\rm Id}_\varnothing) = \alpha_i\circ (g \rhd {\rm Id}_\varnothing) \circ (g'_i \rhd {\rm Id}_\varnothing) = \alpha'_i\circ (g'_i \rhd {\rm Id}_\varnothing)$, and ${\rm St}_W^{-1}(f_i)$ is well defined.

For simplicity we assume that $\St_W^{-1}(f_i)$ and $g_i$ are actually defined on all $W$ but are 0 except on~$W_i$, namely we precompose by the projection $W \twoheadrightarrow W_i$, and thus we set $\St_W^{-1}(f) = \sum_i \St_W^{-1}(f_i)$.

It is now easy to check that $\St_W^{-1}$ is the inverse of $\St_W$. For $v_{\vec{\varepsilon}} \in W$, suppose $v_{\vec{\varepsilon}}\in W_{i_0} \subseteq W$ lies in a simple, or decompose it in the $W_i$'s, and write $v_{\vec{\varepsilon}} = X\cdot w_{i_0}$, $X \in \Uq$. One has $g_{i_0}(v_{\vec{\varepsilon}})= X\cdot v_{\vec{\varepsilon}_{i_0}}$ and for $j\neq {i_0},\ g_j(v_{\vec{\varepsilon}})=0$. Thus:
\begin{align*}
\St_W\big(\St_W^{-1}(f)\big)(v_{\vec{\varepsilon}}) &= \St_W\bigg(\sum_i \alpha_i\circ (g_i \rhd {\rm Id}_\varnothing)\bigg)(v_{\vec{\varepsilon}}) = \sum_i \St_W(\alpha_i\circ (g_i \rhd {\rm Id}_\varnothing))(v_{\vec{\varepsilon}})
\\
&\overset{{\rm Proposition~\ref{propStNat}}}{=} \sum_i \St_W(\alpha_i)(g_i(v_{\vec{\varepsilon}}))
 = \St_W(\alpha_{i_0})(X\cdot v_{\vec{\varepsilon}_{i_0}})
 \\&
\overset{{\rm Proposition~\ref{propStOqmorph}}}{=} X\cdot \St_W(\alpha_{i_0})(v_{\vec{\varepsilon}_{i_0}})= X\cdot a_{i_0}= X \cdot f_{i_0}(w_{i_0})\\
& = f_{i_0}(v_{\vec{\varepsilon}})= f(v_{\vec{\varepsilon}}).
\end{align*}
Symmetrically, let $\alpha \in \Hom_{{\rm Sk}_{{\rm TL}}(\Sigma)}(W \rhd \varnothing,\varnothing)$ and $v_{\vec{\varepsilon}} \in W$. Set $f=\St_W(\alpha)\colon V^{\otimes \abss{\partial \alpha}} \to \mathscr{S}(\Sigma)$, in the definition of $\St_W^{-1}(f)$ one has $\alpha_i = \alpha$ and $v_{\vec{\varepsilon}_i}=w_i$, so $g_i$ is the inclusion $W_i\hookrightarrow W$. If $v_{\vec{\varepsilon}} \in W_i \subseteq W$, one has $\St_W^{-1}(\St_W(\alpha))=\sum_i \alpha\circ (g_i\rhd {\rm Id}_\varnothing) = \alpha\circ ({\rm Id}_W \rhd {\rm Id}_\varnothing) = \alpha$.
\end{proof}

\begin{Proposition}\label{propProdsOnSS}
The algebra structure inherited from the internal endomorphism object structure on~$\mathscr{S}(\Sigma)$ coincides with its usual algebra structure. Namely $A_\Sigma =\mathscr{S}(\Sigma)$ as $\Oq$-comodule algebras.
\end{Proposition}

\begin{proof}
Recall that the product on $\mathscr{S}(\Sigma)$ is given by stacking the left tangle $a$ above the right one~$b$, and the product on $A_\Sigma$ is defined by evaluation and composition maps on internal $\Hom$ objects, in~Definition~\ref{algStrOnIntHom}. Graphically, the stated tangles $a$ and $b$, which we see as morphisms $\alpha$ and $\beta$ in~${\rm Sk}_\mathcal{V}(\Sigma)$ with prescribed inputs $v_{\vec{\varepsilon}}$ and $v_{\vec{\eta}}$, map to the morphism $\alpha \circ ({\rm Id}_{V^{\otimes \abss{\partial \alpha}}}\rhd \beta)$ in~${\rm Sk}_\mathcal{V}(\Sigma)$ with prescribed input $v_{\vec{\varepsilon}} \otimes v_{\vec{\eta}}$, which we see as the stated tangle $a$ above $b$:
\[
\begin{array}{ccc}
\begin{tikzpicture}[scale = 1, baseline = 8pt]
\node [rightymissy, minimum width = 1cm] (A) at (1.5,1) {$\alpha$};
\draw (0,0) -- ++(2,0) ;
\draw (0,0) -- ++ (0,1);
\draw (0,0.75) node[ left = -2pt]{\footnotesize $\varepsilon_1$} -- (A);
\draw (0,0.5) node[ above , rotate = 90]{\tiny $\cdots$} .. controls (1,0.5) .. (A);
\draw (0,0.25) node[ left = -2pt]{\footnotesize $\varepsilon_n$} .. controls (1.5,0.25) .. (A);
\end{tikzpicture} , \begin{tikzpicture}[scale = 1, baseline = 8pt]
\node [rightymissy, minimum width = 1cm] (A) at (1.5,1) {$\beta$};
\draw (0,0) -- ++(2,0) ;
\draw (0,0) -- ++ (0,1);
\draw (0,0.75) node[ left = -2pt]{\footnotesize $\eta_1$} -- (A);
\draw (0,0.5) node[ above , rotate = 90]{\tiny $\cdots$} .. controls (1,0.5) .. (A);
\draw (0,0.25) node[ left = -2pt]{\footnotesize $\eta_m$} .. controls (1.5,0.25) .. (A);
\end{tikzpicture}
&
&
\begin{tikzpicture}[baseline = 10pt]
\draw (0,0) -- ++(4,0) ;
\draw (0,0) -- ++ (0,2);
\begin{scope}[yshift = 1cm]
\node [rightymissy, minimum width = 3cm] (A) at (2.5,1) {$\alpha$};
\draw (0,0.75) node[ left = -2pt]{\footnotesize $\varepsilon_1$} -- (A);
\draw (0,0.5) node[ above , rotate = 90]{\tiny $\cdots$} .. controls (1,0.5) .. (A);
\draw (0,0.25) node[ left = -2pt]{\footnotesize $\varepsilon_n$} .. controls (1.5,0.25) .. (A);
\end{scope}
\begin{scope}[xshift = 2cm]
\node [rightymissy, minimum width = 1cm] (A) at (1.5,1) {$\beta$};
\draw (-2,0.75) node[ left = -2pt]{\footnotesize $\eta_1$} -- (A);
\draw (-2,0.5) node[ above , rotate = 90]{\tiny $\cdots$} .. controls (1,0.5) .. (A);
\draw (-2,0.25) node[ left = -2pt]{\footnotesize $\eta_m$} .. controls (1.5,0.25) .. (A);
\end{scope}
\end{tikzpicture} \hspace{1.8cm}
\\
\updownarrow & & \updownarrow \\
\begin{tikzpicture}[scale = 1, baseline = 8pt]
\node [rightymissy, minimum width = 1cm] (A) at (1.5,1) {$\alpha$};
\draw (0,0) -- ++(2,0);
\draw (0,0) -- ++ (0,1);
\draw (0.5,0) node[above left = -2pt]{\tiny 1} .. controls (0.5,1) .. (A);
\draw (1,0) node[above left = -2pt]{\tiny $\cdots$} .. controls (1,0.7) .. (A);
\draw (1.5,0) node[above left = -2pt]{\tiny $n$} -- (A);
\end{tikzpicture},\ v_{\vec{\varepsilon}} , \
\begin{tikzpicture}[scale = 1, baseline = 8pt]
\node [rightymissy, minimum width = 1cm] (A) at (1.5,1) {$\beta$};
\draw (0,0) -- ++(2,0);
\draw (0,0) -- ++ (0,1);
\draw (0.5,0) node[above left = -2pt]{\tiny 1} .. controls (0.5,1) .. (A);
\draw (1,0) node[above left = -2pt]{\tiny $\cdots$} .. controls (1,0.6) .. (A);
\draw (1.5,0) node[above left = -2pt]{\tiny $m$} -- (A);
\end{tikzpicture},\ v_{\vec{\eta}}
& \quad \longmapsto
&
\begin{tikzpicture}[baseline = 10pt]
\draw (0,0) -- ++(4,0) ;
\draw (0,0) -- ++ (0,2);
\begin{scope}[yshift = 1cm]
\node [rightymissy, minimum width = 3cm] (A) at (2.5,1) {$\alpha$};
\draw (0.5,-1) node[above left = -2pt]{\tiny 1} .. controls (0.5,1) .. (A);
\draw (1,-1) node[above left = -2pt]{\tiny $\cdots$} .. controls (1,0.7) .. (A);
\draw (1.5,-1) node[above left = -2pt]{\tiny $n$} .. controls (1.5,0.6) .. (A);
\end{scope}
\begin{scope}[xshift = 2cm]
\node [rightymissy, minimum width = 1cm] (A) at (1.5,1) {$\beta$};
\draw (0.5,0) node[above left = -2pt]{\tiny 1} .. controls (0.5,1) .. (A);
\draw (1,0) node[above left = -2pt]{\tiny $\cdots$} .. controls (1,0.6) .. (A);
\draw (1.5,0) node[above left = -2pt]{\tiny $m$} -- (A);
\end{scope}
\end{tikzpicture} ,\ v_{\vec{\varepsilon}} \otimes v_{\vec{\eta}}.
\end{array}
\]
The evaluation map ${\rm ev}_{\varnothing,\varnothing}\colon \mathscr{S}(\Sigma)\rhd \varnothing\to\varnothing$ is the image under $\St^{-1}$ of ${\rm Id}_{\mathscr{S}(\Sigma)}$. We have not constructed~$\St^{-1}$ on all $\mathcal{E}$ above, but only on ${\rm TL}$, and it extends by cocontinuity in Proposition~ \ref{propIntHomFree}. The comodule~$\mathscr{S}(\Sigma)$ decomposes as simples as $\mathscr{S}(\Sigma) = \bigoplus_{\alpha \in \mathcal{B}^+} \Uq\cdot \alpha$ where $\mathcal{B}^+$ is the set of $\mathfrak{o}$-ordered simple stated tangles with only + states, see \cite[Theorem~4.6(b)]{Cos}. These stated tangles with only + states are simply a way to represent canonically a tangle without state information, and again in the following we will see $\alpha$ as a morphism in $\Hom_{{\rm Sk}_\mathcal{V}(\Sigma)}\big(V^{\otimes \abss{\partial \alpha}},\varnothing\big)$, which is actually $\psi_{hv}^{-1}(\alpha)$.

 We denote by $W_\alpha = \Uq\cdot \alpha$ and $g_\alpha\colon W_\alpha \to V^{\otimes \abss{\partial \alpha}}$ the inclusion mapping $\alpha$ to its states $v_{\overrightarrow{+\cdots +}}$. Then:
 \[
 {\rm ev}_{\varnothing,\varnothing} =\St^{-1}({\rm Id}_{\mathscr{S}(\Sigma)}) = \oplus_{\alpha \in \mathcal{B}^+} \St^{-1}(W_\alpha \hookrightarrow \mathscr{S}(\Sigma)) = \oplus_{\alpha \in \mathcal{B}^+}\alpha \circ (g_\alpha \rhd {\rm Id}_\varnothing).
 \]
The composition map $c \colon \mathscr{S}(\Sigma) \otimes \mathscr{S}(\Sigma) \to \mathscr{S}(\Sigma)$ is the image under $St$ of the morphism ${\rm ev}_{\varnothing,\varnothing} \circ ({\rm Id}_{\mathscr{S}(\Sigma)} \rhd {\rm ev}_{\varnothing,\varnothing}) \colon (\mathscr{S}(\Sigma) \otimes \mathscr{S}(\Sigma))\rhd \varnothing \to \mathscr{S}(\Sigma) \rhd \varnothing\to\varnothing$. This morphism is the double sum:
\begin{gather*}
\oplus_{\alpha \in \mathcal{B}^+}\oplus_{\beta \in \mathcal{B}^+}(\alpha \circ (g_\alpha \rhd {\rm Id}_\varnothing))\circ ({\rm Id}_{\mathscr{S}(\Sigma)} \rhd (\beta \circ (g_\beta \rhd {\rm Id}_\varnothing)))
\\ \qquad
{}=\oplus_{\alpha \in \mathcal{B}^+}\oplus_{\beta \in \mathcal{B}^+}\alpha \circ ({\rm Id}_{V^{\otimes \abss{\partial \alpha}}}\rhd \beta) \circ (g_\alpha \rhd g_\beta \rhd {\rm Id}_\varnothing) .
\end{gather*}
The product is obtained by applying $\St$ to this morphism. For $a,b \in \mathscr{S}(\Sigma)$, write $a=X\cdot \alpha$ and $b=Y\cdot \beta$ with $X,Y \in \Uq$ and $\alpha,\beta \in \mathcal{B}^+$. Thus $a$ has states $v_{\vec{\varepsilon}} = X \cdot v_{\overrightarrow{+\cdots +}}$ and $b$ has states $v_{\vec{\eta}}=Y \cdot v_{\overrightarrow{+\cdots +}}$. By naturality:
\begin{align*}
c(a\otimes b) &:= \St_{\mathscr{S}(\Sigma)\otimes \mathscr{S}(\Sigma)}({\rm ev}_{\varnothing,\varnothing} \circ ({\rm Id}_{\mathscr{S}(\Sigma)} \rhd {\rm ev}_{\varnothing,\varnothing}))(a\otimes b)
\\
&= \oplus_{\alpha' \in \mathcal{B}^+}\oplus_{\beta' \in \mathcal{B}^+} \St(\alpha' \circ ({\rm Id}_{V^{\otimes \abss{\partial \alpha'}}}\rhd \beta'))((g_{\alpha'} \otimes g_{\beta'})(a \otimes b))
\\
& = \St(\alpha \circ ({\rm Id}_{V^{\otimes \abss{\partial \alpha}}}\rhd \beta))((g_\alpha \otimes g_\beta)(a \otimes b)) = \St(\alpha \circ ({\rm Id}_{V^{\otimes \abss{\partial \alpha}}}\rhd \beta))(v_{\vec{\varepsilon}} \otimes v_{\vec{\eta}})
\end{align*}
which is precisely the usual product of $a$ and $b$ in $\mathscr{S}(\Sigma)$.
\end{proof}
